% La política de privacidad, para imprimir y para el expediente del abogado.
% Dice LO MISMO que /privacidad (server/privacidad.html), vigente desde el
% 11/11/2026: si una cambia, cambia la otra (y el consentimiento con ellas).
\documentclass[10.5pt,a4paper]{article}
% Preámbulo común de los documentos legales (legal/*.tex).
% Se compila cada uno con:   cd legal && latexmk -pdf <archivo>.tex
% La paleta es la del producto y la del manual: si cambia una, cambian las dos.
\usepackage[T1]{fontenc}
\usepackage[utf8]{inputenc}
\usepackage[spanish,es-noshorthands,es-tabla]{babel}
\usepackage[scaled=0.92]{helvet}
\renewcommand{\familydefault}{\sfdefault}
\usepackage{amssymb}
% Los que usan los textos y pdflatex no trae
\DeclareUnicodeCharacter{2713}{\checkmark}
\DeclareUnicodeCharacter{22EF}{\ensuremath{\cdots}}
\DeclareUnicodeCharacter{2192}{\ensuremath{\rightarrow}}
\usepackage[a4paper,top=2cm,bottom=2.2cm,left=2.2cm,right=2.2cm,headheight=14pt]{geometry}
\usepackage[table]{xcolor}
\usepackage{microtype}
\usepackage{enumitem}
\usepackage{booktabs}
\usepackage{tabularx}
\usepackage{array}
\usepackage{longtable}
\usepackage{fancyhdr}
\usepackage{titlesec}
\usepackage{lastpage}
\usepackage[most]{tcolorbox}
\usepackage[hidelinks]{hyperref}

\definecolor{tinta}{HTML}{0B2239}
\definecolor{tinta2}{HTML}{2A4864}
\definecolor{tenue}{HTML}{496A86}
\definecolor{marca}{HTML}{1668C4}
\definecolor{ambar}{HTML}{A67300}
\definecolor{rojo}{HTML}{C2001D}
\definecolor{fondo}{HTML}{F4F7FB}
\definecolor{linea}{HTML}{C3D4E4}
\definecolor{ambarfondo}{HTML}{FDF6E4}
\definecolor{ambarborde}{HTML}{E4C88A}
\definecolor{azulfondo}{HTML}{E7F0FB}
\definecolor{azulborde}{HTML}{C9DDF4}

\color{tinta}
\setlength{\parindent}{0pt}
\setlength{\parskip}{5pt plus 2pt}
\linespread{1.08}
\setlist{itemsep=2pt,topsep=3pt}

\titleformat{\section}{\color{tinta}\large\bfseries}{}{0em}{}
\titleformat{\subsection}{\color{tinta2}\normalsize\bfseries}{}{0em}{}
\titlespacing*{\section}{0pt}{12pt}{4pt}
\titlespacing*{\subsection}{0pt}{8pt}{2pt}

\renewcommand{\arraystretch}{1.25}
\arrayrulecolor{linea}

% Lo que hay que completar a mano: una línea para escribir
\newcommand{\llenar}[1][6cm]{\rule[-2pt]{#1}{0.5pt}}
% Lo que tiene que decidir o completar el abogado o las partes
\newcommand{\completar}[1]{{\color{ambar}[#1]}}
% Una casilla para marcar
\newcommand{\casilla}{\fbox{\rule{0pt}{7pt}\rule{7pt}{0pt}}}

\newtcolorbox{nota}[1][]{colback=azulfondo,colframe=azulborde,boxrule=0.6pt,arc=3pt,
  left=8pt,right=8pt,top=5pt,bottom=5pt,fontupper=\small\color{tinta2},#1}
\newtcolorbox{borrador}{colback=ambarfondo,colframe=ambarborde,boxrule=0.6pt,arc=3pt,
  left=8pt,right=8pt,top=5pt,bottom=5pt,fontupper=\small}

% Pie: qué documento es, su versión y la página
\newcommand{\documento}[2]{%
  \pagestyle{fancy}\fancyhf{}%
  \renewcommand{\headrulewidth}{0pt}%
  \fancyfoot[L]{\scriptsize\color{tenue}#1 · #2}%
  \fancyfoot[R]{\scriptsize\color{tenue}Página \thepage\ de \pageref*{LastPage}}%
}

% Firma: una línea y lo que va debajo
\newcommand{\firma}[1]{\begin{minipage}[t]{0.44\textwidth}\vspace{1.6cm}\rule{\linewidth}{0.5pt}\\{\small\color{tinta2}#1}\end{minipage}}

\begin{document}
\documento{Política de privacidad — Control de ruta}{vigente desde el 11/11/2026}

{\LARGE\bfseries Política de privacidad — Control de ruta}\par
{\small\color{tenue}Vigente desde el 11 de noviembre de 2026 · publicada en \texttt{/privacidad} del servidor}

\textbf{Control de ruta} es la aplicación del conductor y el cobrador de una
cooperativa de transporte que contrató este sistema de control de flota. La
usan personas con una cuenta que su propia cooperativa creó: no tiene
registro abierto ni usuarios anónimos.

\section{Qué datos se recopilan}
\begin{itemize}[leftmargin=14pt]
\item \textbf{Ubicación (también en segundo plano).} Mientras el conductor
  declara que está \emph{en ruta}, el teléfono emite su ubicación de forma
  continua ---también con la pantalla apagada o la aplicación en segundo
  plano---. Es la función central de la app: con esa ubicación se dibuja la
  combi en el mapa de su cooperativa y se calcula la separación con las demás.
  \textbf{Fuera del turno no se emite ubicación}: al declarar «fuera de ruta»
  (o al cerrar sesión) el GPS se apaga, y una ausencia prolongada lo apaga
  sola. En la pantalla de inicio de turno, con la aplicación abierta y el
  permiso ya dado, el teléfono \textbf{mira} la ubicación sin enviarla, sólo
  para darse cuenta de si el conductor ya está manejando sin haber salido a
  ruta: si es así le avisa, y si no contesta en un minuto sale a ruta por él
  ---desde ahí emite como en cualquier turno, y la cooperativa ve que la
  salida fue automática---.
\item \textbf{Identidad de la cuenta.} Nombre, alias opcional, vehículo y rol,
  cargados por la cooperativa. Las contraseñas se guardan cifradas de forma
  irreversible.
\item \textbf{Mensajes del chat de la ruta}: texto, notas de voz y fotos que el
  usuario decide mandar, incluida la emergencia (SOS).
\item \textbf{Actividad del turno}: horas en ruta, vueltas al circuito y desvíos
  del trazado, calculados de la ubicación.
\item \textbf{Lugares puntuales del turno}, con su hora: dónde se declaró en
  ruta, ausente o fuera de ruta ---incluido \textbf{dónde terminó el
  turno}---, dónde se cortó y volvió la señal, dónde la combi estuvo parada y
  dónde salió del trazado. No es el recorrido continuo: son esos puntos.
\item \textbf{Marcas automáticas}: el sistema anota por su cuenta señales como
  «GPS sospechoso», «entrada tardía», «salió solo», «fuera de ruta» o «reloj
  adelantado». Son avisos para que una persona de la cooperativa los mire:
  \textbf{el sistema no sanciona ni decide nada por sí solo}, y el conductor
  puede pedir que se revisen y que se le explique cualquiera.
\item \textbf{Estado del teléfono}, mientras está en ruta y junto con la
  ubicación: el nivel de batería, si está cargando, si el sistema le aplica
  ahorro de batería a la aplicación y la versión de la app. Sirve para que la
  central sepa por qué una combi dejó de verse sin tener que llamar al
  conductor. No se guarda: se ve sólo el último valor.
\item \textbf{Grabación de recorrido}: solo si el conductor la inicia (o acepta
  el pedido de su central), se guarda el trazado recorrido para cargarlo como
  recorrido oficial de la ruta.
\end{itemize}

\section{Para qué se usan}
Para operar la flota de la cooperativa: el mapa en vivo de su central
(Despacho), la separación entre combis, el registro de vueltas y horas con el
que la cooperativa liquida el trabajo, el chat de coordinación y la atención
de emergencias. \textbf{No se usan para publicidad, no se venden ni se
comparten con terceros para sus propios fines.}

\section{Quién los ve}
\begin{itemize}[leftmargin=14pt]
\item \textbf{Su cooperativa}: el personal de Despacho y gerencia de la propia
  cooperativa, dentro de su alcance. Una cooperativa nunca ve datos de otra.
\item \textbf{El operador del sistema}, para administrar el servicio y dar soporte.
\item \textbf{Proveedores de infraestructura}, solo lo mínimo para funcionar: el
  servidor se aloja en Railway, y el fondo del mapa puede pedir imágenes de
  mapa a Geoapify (ese pedido lleva la dirección IP del teléfono, como
  cualquier descarga; no lleva la identidad de la cuenta). Las copias de
  respaldo de la base se guardan además en un almacenamiento aparte,
  \textbf{cifradas}: ese proveedor guarda archivos que no puede leer.
\end{itemize}

\section{Cuánto tiempo se guardan}
\begin{longtable}{>{\raggedright}p{0.5\linewidth}>{\raggedright\arraybackslash}p{0.44\linewidth}}
\toprule
\textbf{Dato} & \textbf{Retención}\\\midrule\endhead
Mensajes del chat (texto) & 30 días\\
Fotos y notas de voz & 30 días. Si una cooperativa junta más de 1500 de cada una, las más viejas se borran antes; queda el aviso de que existieron\\
Mensajes que Despacho guardó como evidencia (con el motivo y quién lo hizo) & 365 días, como un SOS; Despacho puede quitarles la marca antes\\
Vueltas y desvíos de ruta & 120 días\\
Turnos (horas trabajadas) & 365 días\\
Emergencias (SOS) & 365 días\\
Registro de actividad: inicios de sesión & 90 días\\
Registro de actividad: lo del recorrido (desvíos, SOS, cambios de combi) & 365 días\\
Registro de actividad: altas, bajas, cambios de clave y correcciones & 2 años\\
Lugares puntuales del turno (inicio y fin, ausencias, cortes de señal, paradas) y marcas automáticas & 120 días\\
Posición en vivo & Solo la última conocida; no se guarda el recorrido continuo (sólo los lugares puntuales de arriba)\\
Estado del teléfono (batería, ahorro de batería) & Solo el último valor, en memoria; no se guarda\\
Copias de respaldo cifradas & Las borra el almacenamiento a los 30 días\\
Grabaciones de recorrido & Hasta 25 por cooperativa; las más viejas se borran solas\\
\bottomrule
\end{longtable}

\section{Seguridad}
Todos los datos viajan cifrados (HTTPS y WebSocket seguro). Las contraseñas se
guardan con hash irreversible. El acceso está separado por cooperativa y por
rol, y las acciones administrativas quedan auditadas.

\section{Sus derechos}
La cuenta la administra su cooperativa: para corregir sus datos, o para darse
de baja y pedir el borrado de lo que no sea registro laboral de la
cooperativa, el primer canal es su central (Despacho). También puede pedirle
una copia de todo lo que el sistema tiene de usted: Despacho la baja en un
archivo, con sus mensajes, fotos, turnos y vueltas. El formulario de
consentimiento que firma con su cooperativa está en \texttt{/consentimiento}.
También puede escribir al contacto del operador: \llenar[7cm].

\section{Cambios}
Si esta política cambia, la fecha de vigencia de arriba cambia con ella, en
esta misma dirección.
\end{document}
